% T8 -- Context fraction at 1,000 and 2,000 users (slide 60).
\begin{table}[!htbp]
\centering\small
\caption{The gain does not come from giving the model a long history}
\label{tab:context-scale}
\begin{tabular}{@{}lrrrrr@{}}
\toprule
\textbf{Past given} & \textbf{Moments} & \textbf{The rule} & \textbf{900 people} & \textbf{1{,}900 people} & \textbf{Difference} \\
\midrule
30\,\% & 5{,}491 & 47.1\,\% & 50.0\,\% \ ($+2.9$) & 51.4\,\% \ ($+4.3$) & $+1.4$ \\
40\,\% & 4{,}703 & 46.8\,\% & 49.9\,\% \ ($+3.1$) & 51.4\,\% \ ($+4.6$) & $+1.5$ \\
50\,\% & 3{,}922 & 46.6\,\% & 50.1\,\% \ ($+3.5$) & 51.5\,\% \ ($+4.9$) & $+1.4$ \\
60\,\% & 3{,}142 & 47.0\,\% & 50.3\,\% \ ($+3.3$) & 51.2\,\% \ ($+4.2$) & $+0.9$ \\
70\,\% & 2{,}359 & 47.3\,\% & 50.5\,\% \ ($+3.2$) & 51.6\,\% \ ($+4.3$) & $+1.1$ \\
80\,\% \textit{(ref.)} & 1{,}570 & 48.2\,\% & 51.0\,\% \ ($+2.8$) & 52.4\,\% \ ($+4.2$) & $+1.4$ \\
90\,\% & 781 & 50.3\,\% & 51.9\,\% \ ($+1.6$) & 53.6\,\% \ ($+3.3$) & $+1.7$ \\
\bottomrule
\end{tabular}

\figtext{%
\figpara{What is being varied.}{Two trained models, one on 900 training people
and one on 1{,}900, both graded on the same 100 held-out people, and each row
hands them a different amount of that person's day as history. ``30\,\%'' means
the model sees the first 30\,\% of the day and answers the rest.}

\figpara{How to read a row.}{\textbf{Moments} is how many graded moments that row
rests on, and it falls as more of the day is given away as history.
\textbf{The rule} is what ``repeat the last antenna'' scores on those moments.
The two model columns give the score, with the distance to the rule in brackets.
\textbf{Difference} is how much the larger training set was worth on that row.}

\figpara{What it shows.}{The 1{,}900-person model is ahead of the rule by 4.3
even when it has seen only 30\,\% of the day, and by 4.2 when it has seen 80\,\%.
So what the model contributes is essentially independent of how much history it
is handed. Whatever it has learned, it is not something that needs a long run-up
to become useful.}

\figpara{And the last column is the point of the table.}{Going from 900 to
1{,}900 training people is worth between 0.9 and 1.7 on every row alike. More
people helps by the same amount whether the model is given a third of somebody's
day or nine tenths of it, which is one more sign that the two levers are
independent. \stwok, with the same 100 held-out people for both sizes.}}

\vspace{3pt}
\begin{minipage}{\textwidth}\footnotesize
Graded over the whole day rather than a cut of it, 7{,}745 moments, the two sizes
score 50.0\,\% and 51.0\,\% against a rule of 46.8\,\%.
\end{minipage}
\figsource{\nbfinal}{train\_cellid\_final.ipynb}{the context-fraction evaluation at the smaller tiers}
\end{table}
